

\chapter{Requisitos e Arquitetura do Sistema}
\label{chapter:architecture}

\begin{introduction}
Este capítulo define os requisitos que guiaram o desenvolvimento e apresenta a arquitetura geral do sistema, descrevendo como os três subsistemas - mecânico, elétrico e de software - se articulam para cumprir esses requisitos.
\end{introduction}

\section{Requisitos do Sistema}
\label{sec:requisitos}

Os requisitos foram definidos a partir das necessidades concretas do utilizador final e das restrições técnicas impostas pelos componentes selecionados.

\textbf{Requisitos funcionais:}
\begin{itemize}
    \item \textbf{RF1 - Movimento bi-axial:} A plataforma deve mover-se de forma controlada em dois eixos horizontais independentes (X e Y).
    \item \textbf{RF2 - Movimento oscilatório de frequência variável:} A plataforma deve executar movimentos sinusoidais nos dois eixos com frequência, amplitude e duração configuráveis em tempo real.
    \item \textbf{RF-Web - Interface Web embarcada:} Os parâmetros de movimento devem ser configuráveis remotamente via interface web servida pelo próprio ESP32 em modo Access Point (SSID: \texttt{MesaSismica}, IP: \texttt{192.168.4.1}), sem necessidade de ligação USB ou computador externo durante o ensaio.
    \item \textbf{RF4 - Suporte de modelos de teste:} A plataforma deve ter dimensões e resistência suficientes para receber modelos reduzidos de estruturas.
\end{itemize}

\textbf{Requisitos não funcionais:}
\begin{itemize}
    \item \textbf{RNF1 - Custo:} Custo total inferior a 500\,\euro{}.
    \item \textbf{RNF2 - Reprodutibilidade:} Todos os ficheiros de design 3D, firmware e scripts devem estar publicamente disponíveis.
    \item \textbf{RNF3 - Acessibilidade técnica:} Montagem realizável com conhecimentos básicos; estrutura fabricada exclusivamente por impressão 3D; componentes disponíveis em fornecedores genéricos (NEMA17, DRV8825, ESP32, CNC Shield, fusos TR8).
\end{itemize}

\section{Visão Geral da Arquitetura}
\label{sec:visao-geral}

O sistema divide-se em três subsistemas interligados: o mecânico converte o movimento dos motores em deslocamento linear da plataforma; o elétrico controla os atuadores e distribui a energia; o software trata do pré-processamento dos dados sísmicos e executa a lógica de controlo em tempo real.

\section{Subsistema Mecânico}
\label{sec:subsistema-mecanico}

O subsistema mecânico converte o movimento rotativo dos motores em deslocamento linear controlado da plataforma, em dois eixos horizontais ortogonais. A estrutura é composta por uma base fixa (fabricada por impressão 3D) e uma plataforma móvel única (também impressa em 3D), sobre a qual são colocados os modelos estruturais a ensaiar e estão fixos os quatro motores.

Quatro colunas de suporte, dotadas de rótulas nas duas extremidades, ligam a base à plataforma. São essas rótulas que viabilizam o design bi-axial: ao contrário de soluções com guias lineares rígidas por eixo, permitem que as colunas se inclinem livremente em qualquer direção, pelo que a plataforma se desloca simultaneamente nos eixos~X e Y sem constrangimentos entre eles.

\section{Subsistema Elétrico e Eletrónico}
\label{sec:subsistema-eletrico}

O ESP32 é a unidade central de controlo. A CNC Shield~\cite{JoyITCNCKit2}, ligada ao ESP32 por fios individuais, providencia a interface para os quatro \textit{drivers} DRV8825, um por motor. A alimentação é bipartida: o ESP32 é alimentado pela porta USB; os motores recebem energia de uma fonte DC externa dedicada de 12\,V. Separar eletricamente o domínio lógico do domínio de potência evita perturbações nos sinais de controlo. O pino Enable (EN, pino 13) é partilhado entre todos os \textit{drivers}: LOW ativa-os; HIGH desativa-os e reduz o consumo em inatividade.

\section{Subsistema de Software e Firmware}
\label{sec:subsistema-software}

O subsistema divide-se em dois componentes. O script Python opera em modo \textit{offline} e converte dados sísmicos em arrays de passos - foi desenvolvido para validação do pipeline e análise das limitações mecânicas (ver Capítulo~\ref{chapter:validation}), mas não constitui o modo de operação da mesa. O firmware em C++ para o ESP32 implementa o modo de operação principal: movimento sinusoidal bi-axial com parâmetros configuráveis em tempo real via interface web, usando a biblioteca AccelStepper~\cite{AccelStepper} em modo \texttt{runSpeed()}.
